\section{Solutions for Chapter~\ref{ch:chain}}

\paragraph{Exercise~\ref{ex:chain-dimension}.}
\phantomsection
\label{sol:ex:chain-dimension}
The recurrence gives
\[
 t_2=11,\ t_3=36,\ t_4=119,\ t_5=393,\ t_6=1298,
 \ t_7=4287,\ t_8=14159.
\]
For odd length the extra term in \eqref{eq:path-dimension} vanishes,
so $\dim\HH_{\mathrm{path},5}=393$. For even length it is four, giving
$\dim\HH_{\mathrm{path},8}=14163$. The eigenvalue $d>1$ strictly dominates the
other five eigenvalues in modulus. The alphabet bound $6^L$ counts
all words, most of which violate at least one fusion constraint.



\paragraph{Exercise~\ref{ex:chain-projector}.}
\phantomsection
\label{sol:ex:chain-projector}
Let $v_f$ be the conjugated $f$th column in the convention of
\eqref{eq:rho-projector}. Unitarity says $v_f^\dagger v_g=\delta_{fg}$.
For $P_S=\sum_{f\in S}v_fv_f^\dagger$, multiplication gives
$P_S^2=P_S=P_S^\dagger$, and its range is the $|S|$-dimensional
span of those columns. With fusion multiplicities, a channel has
an additional orthonormal multiplicity label, and one sums over
all selected copies. It is the subspace, not a single scalar label,
that defines the projection.



\paragraph{Exercise~\ref{ex:chain-orbit}.}
\phantomsection
\label{sol:ex:chain-orbit}
For a basis vector $|x\rangle$ with stabilizer $K\subset G$, form
\[
 v_\chi=\sum_{g\in G}\overline{\chi(g)}\,|gx\rangle.
\]
Every repeated orbit entry is multiplied by
$\sum_{k\in K}\overline{\chi(k)}$. Character orthogonality makes
this zero unless $\chi|_K=1$. In the surviving case, choose one
representative per coset and divide its character sum by
$\sqrt{|G|/|K|}$. Distinct cosets give orthogonal basis vectors,
so the normalized vector is
\[
 |\chi;x\rangle=
 \frac1{\sqrt{|G|/|K|}}
 \sum_{g\in G/K}\overline{\chi(g)}\,|gx\rangle.
\]
It is independent of the coset representatives precisely because
$\chi|_K=1$. Equivalently, the original full-group sum has norm
$\sqrt{|G||K|}$, not $\sqrt{|G|}$. Replacing $g$ by $h^{-1}g$
in the sum shows that $h|\chi;x\rangle=\chi(h)|\chi;x\rangle$.



\paragraph{Exercise~\ref{ex:chain-noninvertible}.}
\phantomsection
\label{sol:ex:chain-noninvertible}
Using self-adjointness and $Y_\rho^2=1+3Y_\rho$,
\[
 (Y_\rho-2)^2=Y_\rho^2-4Y_\rho+4=5-Y_\rho.
\]
Its eigenvalues are $d-2$ and $-d^{-1}-2$, both nonzero. On a unit
vector,
$5-d\leq\|(Y_\rho-2)s\|^2\leq5+d^{-1}$.
An orthogonal projection has only eigenvalues zero and one;
these values plainly do not have that property. The special
extraction formula in the rank-four vacuum module uses one
particular input vector, not a global projection identity.



\paragraph{Exercise~\ref{ex:chain-boundary}.}
\phantomsection
\label{sol:ex:chain-boundary}
For $L$ physical rho objects between fixed external heights $a,b$,
the dimension is $(N_\rho^L)_{ab}$: matrix multiplication sums
all intermediate fusion histories. Summing the diagonal entries
instead identifies the two ends and permits every compatible
starting height. It changes both the state space and the closing
Hamiltonian terms. Even when the infinite bulk theory agrees,
fixed endpoints can select boundary conformal families or carry
edge states absent from a given periodic sector. The equal leading
growth $d^L$ does not equate their finite spectra.


